\documentclass[11pt]{article}
\usepackage[letterpaper,margin=0.85in]{geometry}\usepackage[T1]{fontenc}\usepackage[utf8]{inputenc}\usepackage{amsmath,amssymb,enumitem,booktabs,array,xcolor,fancyhdr,tabularx}
\setlist{itemsep=2pt,topsep=3pt}\pagestyle{fancy}\fancyhf{}\renewcommand{\headrulewidth}{0.3pt}
\fancyhead[L]{\small 4.1 Interpreting Derivatives in Context}\fancyhead[R]{\small Teacher key}\fancyfoot[C]{\small\thepage}
\definecolor{ans}{RGB}{28,125,74}\newcommand{\A}[1]{\textcolor{ans}{\textbf{#1}}}
\begin{document}

\section*{4.1 Interpreting Derivatives in Context --- teacher key}
\textbf{Question:} What does a derivative tell you when the function is a real quantity, like water in a tank?\par
\subsection*{Self-paced lesson}No clock, no trimming, no skipping, no end mark, no stamp. Every tab is required in order; progress saves on the device so a student can stop and finish the next day.\par
\subsection*{Warm-up review (3 easy items)}\begin{itemize}
\item The derivative of $5x^2-3x$ is... \A{10x -- 3.} Power rule on each term: $5x^2$ gives $10x$, and $-3x$ gives $-3$.
\item A line passes through $(1,2)$ and $(3,8)$. Its slope is... \A{3.} Rise over run: $\dfrac{8-2}{3-1}=3$.
\item If $f'(x)<0$ on an interval, then $f$ is... \A{decreasing.} A negative slope means the graph of $f$ goes down.
\end{itemize}
\subsection*{Explore stations}Four hands-on stations (X1, X2, X3, X4). Station answers are in the teacher view of the lesson page (add \texttt{\#key} to its address).\par
\subsection*{Explore: what did you see}\begin{itemize}
\item In Station 1, at the moment the tank was neither filling nor emptying, the rate reading was... \A{0.} A flat tangent has slope 0, and the slope is the rate.
\item In Station 3, when the step $h$ was small, the tangent-line estimate was... \A{very close to the actual value.} For a small step the curve stays close to its tangent line, so the estimate is good.
\end{itemize}
\subsection*{Mid-lesson check (Huddle tab)}\begin{itemize}
\item $V'(2)=8$ gallons per minute. About how many gallons come in during the next half minute? \A{4.} Change is about rate times time: $8\times0.5=4$ gallons.
\item If $C'(a)<0$, then at $x=a$ the cost is... \A{going down.} A negative rate means the quantity is falling at that point.
\end{itemize}
\subsection*{Targets}\begin{itemize}\item I can say what f?(a) means in a story, with units.
\item I can tell if a quantity is rising or falling from the sign of f?.
\item I can use f(a) and f?(a) to estimate a nearby value.\end{itemize}
\subsection*{Predictions}\begin{itemize}
\item A tank holds $V(t)$ gallons after $t$ minutes, and $V'(5)=-3$. At $t=5$ the tank is... \A{losing water.} A negative derivative means the amount is going down.
\item For the same tank, $V(t)$ is in gallons and $t$ is in minutes. The derivative $V'(t)$ is measured in... \A{gallons per minute.} A rate is the amount that changes divided by the time it takes: gallons per minute.
\end{itemize}
\section*{Huddle script (answers in green)}
\subsection*{Definitions}\begin{itemize}
\item \textbf{Derivative as a rate:} f?(a) is the \A{rate} of change of f at x = a. It is the \A{slope} of the tangent line.
\item \textbf{Units:} The units of f? are the units of f \A{per} the units of x.
\item \textbf{Reading the sign:} If f?(a) > 0, then f is \A{increasing}. If f?(a) < 0, then f is \A{decreasing}.
\item \textbf{Estimating:} f(a + h) $\approx$ f(a) + \A{h} $\cdot$ f?(a). It works best when h is \A{small}.
\end{itemize}
\subsection*{Worked example}
A tank holds V(t) = 40 + 12t -- t\textsuperscript{2} gallons, t in minutes. Find V?(3) and say what it means.\\
\begin{itemize}
\item Differentiate: V?(t) = 12 -- \A{2t}
\item Put in t = 3: V?(3) = 12 -- 6 = \A{6}
\item The units are gallons \A{per minute}
\item At t = 3 minutes, the tank is \A{filling} at 6 gallons per minute.
\end{itemize}
\subsection*{Label the picture}
\begin{enumerate}
\item The curve: its height is the \A{amount}.
\item The tangent line: its slope is the \A{rate}.
\item The horizontal axis: the \A{time} or input.
\end{enumerate}
\section*{Every level, with answers}
\subsection*{Say what it means}
\noindent\textbf{Level 1.} $T(t)$ is the temperature, in degrees Fahrenheit, of an oven $t$ minutes after it is turned on. Also $T'(7)=8$. What does this say?\\ (A) The temperature is 8 degrees when t = 7. \quad \A{(B) When t = 7, the temperature is rising at 8 degrees per minute.} \quad (C) When t = 7, the temperature is rising at 8 minutes per degree. \quad (D) When t = 7, the temperature is falling at 8 degrees per minute.\\ \A{Answer: B} \quad The sign is positive, so the temperature is rising. The size is 8 and the units are degrees per minute.\par\medskip
\noindent\textbf{Level 2.} $C(x)$ is the cost, in dollars, of making $x$ shirts. Also $C'(80)=-8$. What does this say?\\ (A) When x = 80, the cost is rising at 8 dollars per shirt. \quad (B) The cost is 8 dollars when x = 80. \quad \A{(C) When x = 80, the cost is falling at 8 dollars per shirt.} \quad (D) When x = 80, the cost is falling at 8 shirts per dollar.\\ \A{Answer: C} \quad The sign is negative, so the cost is falling. The size is 8 and the units are dollars per shirt.\par\medskip
\noindent\textbf{Level 3.} $H(t)$ is the height, in feet, of a tree $t$ years after it was planted. Which equation says: "When t = 1, the height is growing at 7 feet per year"?\\ (A) $H(1)=7$ \quad \A{(B) $H'(1)=7$} \quad (C) $H'(7)=1$ \quad (D) $H(7)=1$\\ \A{Answer: B} \quad A rate is a derivative, so use $H'$. The input is 1 and the rate is 7: $H'(1)=7$.\par\medskip
\noindent\textbf{Level 4.} $H(t)$ is the height, in feet, of a tree $t$ years after it was planted. Which equation says: "When t = 11, the height is growing at 3 feet per year"?\\ (A) $H(11)=3$ \quad \A{(B) $H'(11)=3$} \quad (C) $H(3)=11$ \quad (D) $H'(3)=11$\\ \A{Answer: B} \quad A rate is a derivative, so use $H'$. The input is 11 and the rate is 3: $H'(11)=3$.\par\medskip
\subsection*{Units of a rate}
\noindent\textbf{Level 1.} $T(t)$ is the temperature, in degrees Fahrenheit, of an oven $t$ minutes after it is turned on. What are the units of $T'(t)$?\\ \A{(A) degrees per minute} \quad (B) degrees minutes \quad (C) degrees \quad (D) minutes per degrees\\ \A{Answer: A} \quad The output is in degrees and the input is in minutes, so the derivative is in degrees per minute.\par\medskip
\noindent\textbf{Level 2.} $C(x)$ is the cost, in dollars, of making $x$ shirts. What are the units of $C''(x)$?\\ (A) dollars per shirt \quad \A{(B) dollars per shirt per shirt} \quad (C) dollars per shirt squared times shirt \quad (D) dollars\\ \A{Answer: B} \quad The output is in dollars and the input is in shirts, so the second derivative is in dollars per shirt per shirt.\par\medskip
\noindent\textbf{Level 3.} $W(h)$ is the volume, in cubic feet, of water in a rain barrel after $h$ hours of rain. What are the units of $W'(h)$?\\ (A) hours per cubic feet \quad (B) cubic feet hours \quad (C) cubic feet \quad \A{(D) cubic feet per hour}\\ \A{Answer: D} \quad The output is in cubic feet and the input is in hours, so the derivative is in cubic feet per hour.\par\medskip
\noindent\textbf{Level 4.} $W(h)$ is the volume, in cubic feet, of water in a rain barrel after $h$ hours of rain. What are the units of $W'(h)$?\\ (A) cubic feet \quad (B) hours per cubic feet \quad \A{(C) cubic feet per hour} \quad (D) cubic feet hours\\ \A{Answer: C} \quad The output is in cubic feet and the input is in hours, so the derivative is in cubic feet per hour.\par\medskip
\subsection*{Estimate a nearby value}
\noindent\textbf{Level 1.} $T(t)$ is the temperature, in degrees Fahrenheit, of an oven $t$ minutes after it is turned on. Right now $t=10$, with $T(10)=330$ and $T'(10)=8$ degrees per minute. Estimate $T(13)$.\\ \A{Answer: 354 degrees} \quad The step is 3 minutes. Change $\approx$ 8 $\times$ 3 = 24. So $T(13)\approx330+24=354$.\par\medskip
\noindent\textbf{Level 2.} $C(x)$ is the cost, in dollars, of making $x$ shirts. Right now $x=83$, with $C(83)=790$ and $C'(83)=-2$ dollars per shirt. Estimate $C(80)$.\\ \A{Answer: 796 dollars} \quad The step is -3 shirts. Change $\approx$ -2 $\times$ -3 = 6. So $C(80)\approx790+6=796$.\par\medskip
\noindent\textbf{Level 3.} $H(t)=-t^{3}-3t^{2}+13t+12$ is the height, in feet, of a tree $t$ years after it was planted. Use $H(1)$ and $H'(1)$ to estimate $H(4)$.\\ \A{Answer: 33 feet} \quad $H(1)=21$ and $H'(1)=4$. The estimate is $21+3\cdot(4)=33$.\par\medskip
\noindent\textbf{Level 4.} $H(t)=-3t^{2}+10t+50$ is the height, in feet, of a tree $t$ years after it was planted. Use $H(6)$ and $H'(6)$ to estimate $H(7)$.\\ \A{Answer: -24 feet} \quad $H(6)=2$ and $H'(6)=-26$. The estimate is $2+1\cdot(-26)=-24$.\par\medskip
\subsection*{Find the rate}
\noindent\textbf{Level 1.} $T(t)=-2t^{2}+13t+48$ is the temperature, in degrees Fahrenheit, of an oven $t$ minutes after it is turned on. Find $T'(3)$.\\ \A{Answer: 1 degrees per minute} \quad $T'(t)=-4t+13$. Putting in $t=3$ gives $1$ degrees per minute.\par\medskip
\noindent\textbf{Level 2.} $C(x)=2\sqrt{x}$ is the cost, in dollars, of making $x$ shirts. Find $C'(16)$.\\ \A{Answer: 0.25 dollars per shirt} \quad $C'(x)=\dfrac{2}{2\sqrt{x}}$. Putting in $x=16$ gives $0.25$ dollars per shirt.\par\medskip
\noindent\textbf{Level 3.} $H(t)=70e^{0.2t}$ is the height, in feet, of a tree $t$ years after it was planted. Find $H'(0)$.\\ \A{Answer: 14 feet per year} \quad $H'(t)=70\cdot0.2e^{0.2t}$. Putting in $t=0$ gives $14$ feet per year.\par\medskip
\noindent\textbf{Level 4.} $H(t)=(t^{2}+1)^{3}$ is the height, in feet, of a tree $t$ years after it was planted. Find $H'(1)$.\\ \A{Answer: 24 feet per year} \quad $H'(t)=3(t^2+1)^{2}\cdot2t$. Putting in $t=1$ gives $24$ feet per year.\par\medskip
\section*{Common misconceptions}\begin{tabularx}{\textwidth}{@{}X X@{}}\toprule Students say & Say instead\\\midrule
V?(3) = 6 means the tank holds 6 gallons at 3 minutes. & V(3) is the amount. V?(3) is the rate: 6 gallons per minute.\\
The units of V? are gallons. & A rate has units of the output per units of the input: gallons per minute.\\
A negative derivative means the quantity is negative. & It means the quantity is going down. The quantity itself can still be large.\\
The estimate f(a) + h f?(a) is exact. & It is an estimate. It is best for small h and gets worse as h grows.\\
Forgetting the units or the direction when writing what f? means. & Write rising or falling, then the size, then "units per unit".\\
\bottomrule\end{tabularx}
\end{document}